\documentclass[a4paper,11pt]{article}
\usepackage[utf8]{inputenc}
\usepackage{amsmath, amssymb}
\usepackage[top=1cm, bottom=1cm, left=1cm, right=1cm]{geometry}
\renewcommand{\arraystretch}{1.15}
\setlength{\parindent}{0pt}

\begin{document}
\pagestyle{empty}

% ===================== FILA SUPERIOR: A y B =====================
\noindent
\begin{minipage}[t]{0.48\textwidth}
    \textbf{Evaluación A} \hfill \textbf{Nombre y Apellido: \underline{\hspace{4cm}}}

    \vspace{0.2cm}
    \begin{enumerate}
        \item \textbf{(2 pts)} Un estanque abierto tiene nivel libre $H=1{,}8\,\mathrm{m}$ por encima de un orificio pequeño en el fondo. Calcule la velocidad de salida del agua.
        \item \textbf{(2 pts)} Agua incompresible fluye horizontalmente desde una tubería de diámetro $D_1=80\,\mathrm{mm}$ a otra de $D_2=40\,\mathrm{mm}$. Si en la primera sección la velocidad es $v_1=1{,}5\,\mathrm{m/s}$, determine la velocidad en la segunda sección y la diferencia de presiones entre ambas.
        \item \textbf{(2 pts)} Un venturí de agua tiene $D_\text{entrada}=10\,\mathrm{cm}$ y garganta $D_t=5\,\mathrm{cm}$. Si a la entrada $v_1=2{,}0\,\mathrm{m/s}$, calcule la velocidad en la garganta y la diferencia de presión con respecto a la entrada.
        \item \textbf{(2 pts)} Agua a $20^\circ\mathrm{C}$ con $\nu=1{,}0\times10^{-6}\,\mathrm{m^2/s}$ circula en una cañería de $D=1{,}0\,\mathrm{in}$ con velocidad media $V=1{,}2\,\mathrm{m/s}$. Calcule el número de Reynolds \(\displaystyle Re=\frac{VD}{\nu}\) e indique el régimen.
        \item \textbf{(2 pts)} En un canal de ancho $b=0{,}80\,\mathrm{m}$ y tirante $h=0{,}20\,\mathrm{m}$ circula agua con $V=0{,}60\,\mathrm{m/s}$ y $\nu=1{,}2\times10^{-6}\,\mathrm{m^2/s}$. Use el diámetro hidráulico \(D_h=\dfrac{4A}{P}\) con \(A=bh\) y \(P=b+2h\) para calcular \(Re=\dfrac{V D_h}{\nu}\) y clasificar el régimen.
    \end{enumerate}
\end{minipage}\hfill
\begin{minipage}[t]{0.48\textwidth}
    \textbf{Evaluación B} \hfill \textbf{Nombre y Apellido: \underline{\hspace{4cm}}}

    \vspace{0.2cm}
    \begin{enumerate}
        \item \textbf{(2 pts)} Un depósito abierto tiene el nivel libre $H=3{,}5\,\mathrm{m}$ por encima de un orificio circular de diámetro $d=12\,\mathrm{mm}$. Calcule la velocidad de salida y el caudal.
        \item \textbf{(2 pts)} En una corriente de agua, un tubo de Pitot registra una altura de estancamiento $h=0{,}25\,\mathrm{m}$. Estime la velocidad de la corriente.
        \item \textbf{(2 pts)} Un aceite con $\nu=30\,\mathrm{cSt}$ escurre en un tubo de $D=50\,\mathrm{mm}$ con $V=0{,}50\,\mathrm{m/s}$. (Use $1\,\mathrm{cSt}=10^{-6}\,\mathrm{m^2/s}$). Calcule \(Re=\dfrac{VD}{\nu}\) y determine el régimen.
        \item \textbf{(2 pts)} Agua con $\nu=1{,}0\times10^{-6}\,\mathrm{m^2/s}$ en un conducto de $D=20\,\mathrm{mm}$. Determine la \emph{máxima} velocidad media para que el flujo sea laminar (\(Re<2000\)): \(V_{\max}=\dfrac{Re_\text{crit}\,\nu}{D}\).
        \item \textbf{(2 pts)} Un depósito grande descarga por una tobera situada $z=2{,}0\,\mathrm{m}$ por debajo del nivel libre. Desprecie pérdidas y determine la velocidad de salida.
    \end{enumerate}
\end{minipage}

\vspace{0.6cm}

% ===================== FILA INFERIOR: C y D =====================
\noindent
\begin{minipage}[t]{0.48\textwidth}
    \textbf{Evaluación C} \hfill \textbf{Nombre y Apellido: \underline{\hspace{4cm}}}

    \vspace{0.2cm}
    \begin{enumerate}
        \item \textbf{(2 pts)} Un estanque abierto tiene nivel libre $H=2{,}2\,\mathrm{m}$ por encima de un orificio pequeño en el fondo. Calcule la velocidad de salida del agua.
        \item \textbf{(2 pts)} Agua fluye horizontalmente desde $D_1=60\,\mathrm{mm}$ a $D_2=30\,\mathrm{mm}$. Si $v_1=1{,}8\,\mathrm{m/s}$, determine $v_2$ y la diferencia de presiones entre ambas secciones.
        \item \textbf{(2 pts)} Venturí: $D_\text{entrada}=8\,\mathrm{cm}$, $D_t=4\,\mathrm{cm}$, $v_1=1{,}5\,\mathrm{m/s}$. Calcule la velocidad en la garganta y la diferencia de presión (ideal).
        \item \textbf{(2 pts)} Agua con $\nu=1{,}0\times10^{-6}\,\mathrm{m^2/s}$ en $D=25\,\mathrm{mm}$ y $V=0{,}90\,\mathrm{m/s}$. Calcule \(Re=\dfrac{VD}{\nu}\) y clasifique el régimen.
        \item \textbf{(2 pts)} Canal: $b=1{,}00\,\mathrm{m}$, $h=0{,}15\,\mathrm{m}$, $V=0{,}75\,\mathrm{m/s}$, $\nu=1{,}1\times10^{-6}\,\mathrm{m^2/s}$. Use \(D_h=\dfrac{4A}{P}\) para hallar \(Re=\dfrac{V D_h}{\nu}\) y el régimen.
    \end{enumerate}
\end{minipage}\hfill
\begin{minipage}[t]{0.48\textwidth}
    \textbf{Evaluación D} \hfill \textbf{Nombre y Apellido: \underline{\hspace{4cm}}}

    \vspace{0.2cm}
    \begin{enumerate}
        \item \textbf{(2 pts)} Depósito abierto con nivel libre $H=3{,}0\,\mathrm{m}$ sobre un orificio de diámetro $d=10\,\mathrm{mm}$. Calcule la velocidad de salida y el caudal.
        \item \textbf{(2 pts)} Un Pitot registra $h=0{,}20\,\mathrm{m}$ en agua. Estime la velocidad de la corriente.
        \item \textbf{(2 pts)} Aceite con $\nu=20\,\mathrm{cSt}$ en tubo $D=40\,\mathrm{mm}$, $V=0{,}60\,\mathrm{m/s}$. Convierta $\nu$ y calcule \(Re=\dfrac{VD}{\nu}\). Clasifique el régimen.
        \item \textbf{(2 pts)} Agua con $\nu=1{,}0\times10^{-6}\,\mathrm{m^2/s}$ en $D=32\,\mathrm{mm}$. Determine $V_{\max}$ para que $Re<2000$.
        \item \textbf{(2 pts)} Un depósito descarga por una tobera $z=1{,}5\,\mathrm{m}$ por debajo del nivel libre. Desprecie pérdidas y determine la velocidad de salida.
    \end{enumerate}
\end{minipage}

\end{document}
